\chapter{Lezione 18 --- Controllo ottimo e programmazione dinamica}

\section{Introduzione al controllo ottimo}

Vogliamo impostare un problema di controllo in retroazione che non si limiti a
stabilizzare il sistema o a imporre certi poli, ma che soddisfi un
\textbf{criterio di ottimalità}.

L'idea è minimizzare una funzione di costo che tenga conto contemporaneamente di:
\begin{itemize}
    \item andamento della traiettoria di stato;
    \item sforzo di controllo;
    \item eventuali vincoli o penalità sullo stato finale.
\end{itemize}

Il controllo ottimo è particolarmente importante per sistemi MIMO e per problemi
in cui non basta ``chiudere bene il sistema'', ma bisogna anche farlo nel modo
più conveniente possibile.

Un'applicazione moderna di questa idea è il \textbf{Model Predictive Control} (MPC),
che può essere visto come uno sviluppo del controllo ottimo.

\section{Due approcci al controllo ottimo}

Negli appunti vengono richiamati due approcci classici:
\begin{enumerate}
    \item \textbf{Calcolo delle variazioni / principio di Pontryagin};
    \item \textbf{Programmazione dinamica / principio di Bellman}.
\end{enumerate}

Entrambi descrivono lo stesso problema di ottimo, ma con strumenti diversi.
In questa lezione ci concentriamo sulla \textbf{programmazione dinamica} per sistemi
a tempo discreto.

\section{Principio di Bellman e idea del cost-to-go}

Per introdurre il metodo si considera il classico esempio del percorso ottimo:
si deve andare da un punto iniziale \(x_A\) a un punto finale \(x_B\), attraversando
eventuali punti di transito organizzati per \emph{stadi}.

Ogni arco del grafo ha un costo, che può rappresentare:
\begin{itemize}
    \item tempo di percorrenza;
    \item consumo di carburante;
    \item rischio;
    \item una qualsiasi grandezza da minimizzare.
\end{itemize}

Il \textbf{principio di Bellman} dice che:

\begin{quote}
Se un percorso è ottimo, allora ogni suo tratto finale è a sua volta ottimo
rispetto al punto da cui parte.
\end{quote}

Questo porta al concetto di \textbf{cost-to-go}:

\begin{quote}
il cost-to-go è il costo minimo necessario per andare dallo stato attuale fino
alla destinazione finale.
\end{quote}

La strategia consiste quindi nel lavorare \emph{all'indietro}, dallo stadio finale
verso lo stadio iniziale.

\section{Problema dinamico discreto generale}

Consideriamo un sistema discreto generico
\[
x(i+1)=f(x(i),u(i)),
\]
e una funzione di costo additiva del tipo
\[
J=\sum_{i=0}^{N-1} h_i(x(i),u(i)).
\]

Il problema è:

\begin{quote}
determinare la sequenza di controlli
\[
u(0),u(1),\dots,u(N-1)
\]
che minimizza il costo \(J\), compatibilmente con la dinamica del sistema.
\end{quote}

\subsection{Definizione di cost-to-go ottimo}

Si definisce il costo ottimo a partire dallo stadio \(i\) come
\[
J_i^o(x(i))
=
\min_{u(i),u(i+1),\dots,u(N-1)}
\sum_{k=i}^{N-1} h_k(x(k),u(k)).
\]

In particolare:
\[
J_0^o(x(0))=\min J.
\]

\subsection{Equazione di Bellman}

Poiché il costo è additivo, vale la ricorsione
\[
J_i^o(x(i))
=
\min_{u(i)}
\left[
h_i(x(i),u(i))
+
J_{i+1}^o\bigl(f(x(i),u(i))\bigr)
\right].
\]

Questa è la \textbf{equazione di Bellman}.

L'inizializzazione è data dalla condizione terminale. Nel caso più semplice:
\[
J_N^o(x(N))=0,
\]
oppure, più in generale, si può assegnare una penalità terminale.

\subsection{Osservazione}

In generale non esiste una soluzione analitica esplicita del problema.
La strategia ottima dipende dallo stato in cui il sistema si trova, quindi il controllo
ottimo è in generale una legge di retroazione
\[
u(i)=\mu_i(x(i)).
\]

\section{Caso particolare: controllo ottimo LQ}

Un caso molto importante, e in effetti uno dei pochi con soluzione analitica,
è il problema \textbf{LQ}:
\begin{itemize}
    \item \textbf{L} = sistema lineare;
    \item \textbf{Q} = costo quadratico.
\end{itemize}

\subsection{Sistema}

Consideriamo il sistema lineare discreto
\[
x_{k+1}=Ax_k+Bu_k,
\qquad
x_0=\hat x,
\]
dove
\[
x_k\in\mathbb{R}^n,\qquad
u_k\in\mathbb{R}^m.
\]

\subsection{Indice di costo}

Il costo è scelto nella forma
\[
J
=
\sum_{k=0}^{N-1}
\left(
x_k^T V x_k + u_k^T P u_k
\right)
+
x_N^T V_N x_N,
\]
dove:
\begin{itemize}
    \item \(V=V^T\succeq 0\): penalità sullo stato;
    \item \(P=P^T\succ 0\): penalità sul controllo;
    \item \(V_N=V_N^T\succeq 0\): penalità sullo stato finale.
\end{itemize}

L'ultimo termine penalizza lo stato di arrivo al tempo \(N\).

\section{Forma attesa della soluzione}

Cerchiamo una legge ottima del tipo
\[
u_k^o=-H_k x_k,
\]
cioè una retroazione lineare sullo stato.

La cosa non è banale a priori, ma la programmazione dinamica mostrerà che,
nel caso LQ, questa forma è effettivamente ottima.

\section{Richiamo utile: minimo di una forma quadratica}

Negli appunti compare il seguente fatto elementare.

Se
\[
Q\succ 0,
\]
allora la funzione
\[
\phi(z)=z^TQz+2c^Tz
\]
ha minimo unico nel punto
\[
z^\star=-Q^{-1}c.
\]

Questo risultato verrà usato ripetutamente nella minimizzazione rispetto a \(u_k\).

\section{Passo finale della ricorsione: stadio \(N-1\)}

Partiamo dallo stadio finale utile, cioè \(k=N-1\).

Poiché
\[
x_N=Ax_{N-1}+Bu_{N-1},
\]
si ha
\[
J_{N-1}^o(x_{N-1})
=
\min_{u_{N-1}}
\left[
x_{N-1}^T V x_{N-1}
+
u_{N-1}^T P u_{N-1}
+
x_N^T V_N x_N
\right].
\]

Sostituendo \(x_N=Ax_{N-1}+Bu_{N-1}\):
\[
\begin{aligned}
J_{N-1}^o(x_{N-1})
&=
\min_{u_{N-1}}
\Big[
x_{N-1}^T V x_{N-1}
+
u_{N-1}^T P u_{N-1}\\
&\qquad\qquad
+
(Ax_{N-1}+Bu_{N-1})^T
V_N
(Ax_{N-1}+Bu_{N-1})
\Big].
\end{aligned}
\]

Sviluppando:
\[
\begin{aligned}
J_{N-1}^o(x_{N-1})
&=
x_{N-1}^T(V+A^TV_NA)x_{N-1}\\
&\quad+
\min_{u_{N-1}}
\Big[
u_{N-1}^T(P+B^TV_NB)u_{N-1}
+
2x_{N-1}^T A^TV_NBu_{N-1}
\Big].
\end{aligned}
\]

Ponendo
\[
Q=P+B^TV_NB,
\qquad
c=B^TV_NA\,x_{N-1},
\]
e usando il risultato precedente, si ottiene
\[
u_{N-1}^o(x_{N-1})
=
-(P+B^TV_NB)^{-1}B^TV_NA\,x_{N-1}.
\]

Definendo
\[
H_{N-1}
=
(P+B^TV_NB)^{-1}B^TV_NA,
\]
si ha
\[
u_{N-1}^o=-H_{N-1}x_{N-1}.
\]

Sostituendo il valore ottimo, il costo ottimo allo stadio \(N-1\) assume forma quadratica:
\[
J_{N-1}^o(x_{N-1})=x_{N-1}^T T_{N-1}x_{N-1},
\]
con
\[
T_{N-1}
=
V+A^TV_NA
-
A^TV_NB(P+B^TV_NB)^{-1}B^TV_NA.
\]

\section{Passo ricorsivo generale}

Supponiamo ora che, a uno stadio generico \(k+1\), il cost-to-go ottimo sia già noto
nella forma quadratica
\[
J_{k+1}^o(x_{k+1})=x_{k+1}^T T_{k+1}x_{k+1}.
\]

Allora allo stadio \(k\):
\[
J_k^o(x_k)
=
\min_{u_k}
\left[
x_k^TVx_k+u_k^TPu_k+J_{k+1}^o(x_{k+1})
\right],
\]
con
\[
x_{k+1}=Ax_k+Bu_k.
\]

Sostituendo:
\[
\begin{aligned}
J_k^o(x_k)
&=
\min_{u_k}
\Big[
x_k^TVx_k
+
u_k^TPu_k
+
(Ax_k+Bu_k)^TT_{k+1}(Ax_k+Bu_k)
\Big].
\end{aligned}
\]

Sviluppando:
\[
\begin{aligned}
J_k^o(x_k)
&=
x_k^T(V+A^TT_{k+1}A)x_k\\
&\quad+
\min_{u_k}
\Big[
u_k^T(P+B^TT_{k+1}B)u_k
+
2x_k^TA^TT_{k+1}Bu_k
\Big].
\end{aligned}
\]

Usando di nuovo il minimo della forma quadratica:
\[
u_k^o(x_k)
=
-(P+B^TT_{k+1}B)^{-1}B^TT_{k+1}A\,x_k.
\]

Definendo
\[
H_k
=
(P+B^TT_{k+1}B)^{-1}B^TT_{k+1}A,
\]
si ottiene la legge ottima
\[
u_k^o=-H_kx_k.
\]

Sostituendo nel costo:
\[
J_k^o(x_k)=x_k^TT_kx_k,
\]
dove
\[
T_k
=
V
+
A^TT_{k+1}A
-
A^TT_{k+1}B
(P+B^TT_{k+1}B)^{-1}
B^TT_{k+1}A.
\]

\section{Ricorsione di Riccati discreta}

La ricorsione
\[
T_k
=
V
+
A^TT_{k+1}A
-
A^TT_{k+1}B
(P+B^TT_{k+1}B)^{-1}
B^TT_{k+1}A
\]
è l'\textbf{equazione di Riccati a tempo discreto} (all'indietro).

L'inizializzazione è
\[
T_N=V_N.
\]

Quindi l'algoritmo di soluzione del problema LQ su orizzonte finito è:

\begin{enumerate}
    \item fissare la condizione terminale
    \[
    T_N=V_N;
    \]
    \item calcolare all'indietro
    \[
    T_{N-1},T_{N-2},\dots,T_0;
    \]
    \item ad ogni passo ricavare
    \[
    H_k=(P+B^TT_{k+1}B)^{-1}B^TT_{k+1}A;
    \]
    \item applicare il controllo ottimo
    \[
    u_k^o=-H_kx_k.
    \]
\end{enumerate}

\section{Forma finale della soluzione}

Riassumendo, la soluzione del problema LQ discreto a orizzonte finito è:
\[
u_k^o=-H_kx_k,
\]
con
\[
H_k=(P+B^TT_{k+1}B)^{-1}B^TT_{k+1}A,
\]
e
\[
J_k^o(x_k)=x_k^TT_kx_k,
\]
dove \(T_k\) soddisfa la ricorsione di Riccati
\[
T_k
=
V
+
A^TT_{k+1}A
-
A^TT_{k+1}B
(P+B^TT_{k+1}B)^{-1}
B^TT_{k+1}A,
\qquad
T_N=V_N.
\]

\section{Osservazioni importanti}

\subsection{Controllo in retroazione}

La soluzione ottima non è una sequenza aperta di controlli, ma una \textbf{retroazione}
lineare sullo stato:
\[
u_k^o=-H_kx_k.
\]

Questa è una delle ragioni del successo del metodo:
il controllo ottimo è robusto rispetto a perturbazioni e incertezze sullo stato,
perché viene ricalcolato a ogni istante a partire da \(x_k\).

\subsection{Guadagno tempo-variante}

Anche se il sistema di partenza è tempo-invariante, il guadagno ottimo
\[
H_k
\]
dipende in generale da \(k\), perché dipende da \(T_{k+1}\) e quindi
dalla distanza dal tempo finale \(N\).

Dunque, su orizzonte finito, il controllo ottimo è in generale \textbf{tempo-variante}.

\subsection{Stabilità}

Negli appunti compare giustamente la domanda:

\begin{quote}
\emph{Ma siamo sicuri che il sistema sia stabile?}
\end{quote}

A questo stadio della trattazione la risposta non è ancora automatica:
la ricorsione fornisce il controllo ottimo su orizzonte finito, ma la stabilità
asintotica del sistema chiuso richiede un'analisi ulteriore.

Questo porterà naturalmente, nelle lezioni successive, al caso a orizzonte infinito,
in cui il guadagno tende a una matrice costante e compare l'equazione algebrica di Riccati.

\section{Interpretazione dei pesi \(V\), \(P\), \(V_N\)}

I tre pesi del problema hanno significato fisico chiaro:
\begin{itemize}
    \item \(V\): penalizza deviazioni dello stato lungo tutta la traiettoria;
    \item \(P\): penalizza lo sforzo di controllo;
    \item \(V_N\): penalizza lo stato finale.
\end{itemize}

Se aumento \(P\), il controllo ottimo diventa più prudente:
si cerca di usare meno attuazione.

Se aumento \(V\) o \(V_N\), si privilegia il rispetto della traiettoria e/o del target finale,
accettando anche uno sforzo di controllo maggiore.

\section{Sintesi della lezione}

I risultati principali della lezione sono:

\begin{itemize}
    \item il controllo ottimo si può affrontare con il principio di Bellman;
    \item si definisce il cost-to-go ottimo
    \[
    J_i^o(x(i));
    \]
    \item nel caso lineare-quadratico discreto:
    \[
    x_{k+1}=Ax_k+Bu_k,
    \qquad
    J=\sum_{k=0}^{N-1}(x_k^TVx_k+u_k^TPu_k)+x_N^TV_Nx_N;
    \]
    \item il cost-to-go ottimo è quadratico:
    \[
    J_k^o(x_k)=x_k^TT_kx_k;
    \]
    \item il controllo ottimo è lineare nello stato:
    \[
    u_k^o=-H_kx_k;
    \]
    \item il guadagno ottimo è
    \[
    H_k=(P+B^TT_{k+1}B)^{-1}B^TT_{k+1}A;
    \]
    \item la matrice \(T_k\) soddisfa la ricorsione di Riccati all'indietro:
    \[
    T_k
    =
    V
    +
    A^TT_{k+1}A
    -
    A^TT_{k+1}B
    (P+B^TT_{k+1}B)^{-1}
    B^TT_{k+1}A,
    \qquad
    T_N=V_N.
    \]
\end{itemize}